\section{Team Structure and Runtime Interventions (RQ2)}
\label{sec:structure}

On tasks that one context can hold, one strong owner matched or beat every team we built, at a
fraction of the cost. Extra agents won on the tasks they were designed on and tied or lost on tasks
they were not, usually at two to five times the cost.

\para{Small tasks are not delegated, big ones are} On the first suite (28 September) the default
team never delegated, and a single agent was as good and slightly cheaper (10.9 against 12.1
minutes, 1.24M against 1.41M tokens). On bigger multi-module tasks (29 September) the lead
delegated in all six runs of two tasks and in none of the research tasks; all 36 runs were solved,
so structure could only change time and cost.

\para{Structures built for hard tasks did not transfer} From 8 October we built one structure a day
for the hard tasks the shipped team missed (\cref{tab:structures}). Diagnoses of those misses found
that each was a single wrong reading of an ambiguous or silent sentence in the request, made the
same way in most runs, so votes and extra reviewers had little to work with. Reading rules on every
node even made results worse out of sample, because the rules licensed the model to skip whatever
the text did not state.

\begin{table}[t]
\centering
\footnotesize
\caption{Team structures built in V1 for hard tasks, each against its same-batch base (arm against
base, runs passed). Design tasks are the tasks a structure was built for; held-out tasks were not
used to build it.}
\label{tab:structures}
\begin{tabularx}{\textwidth}{@{}>{\raggedright\arraybackslash}p{36mm} L L >{\raggedright\arraybackslash}p{19mm} >{\raggedright\arraybackslash}p{27mm}@{}}
\toprule
Structure & Design tasks & Held-out tasks & Cost per run & Decision \\
\midrule
Reading rules on every node & 9/15 against 4/15 & 2/10 against 4/10 & about \$1.00 against \$0.76 &
  Rejected: the rules shrank the work \\
Three builders and an evidence judge (ensemble) & 15/24 against 4/15; the first solves of one task &
  3/10 against 8/20 & \$2.85 against \$0.53 & Rejected \\
Three readers and a reconciler & 4/10 against 4/21 & 1/9 against 11/28 & \$4.01 against \$1.31 &
  Rejected \\
Acceptance checker started beside the build & 3/5 against 2/15 (detail-heavy tasks) & 8/10 against
  7/10 (fresh set) & \$2.03 against \$1.58 & Adopted as a gate, only for specification-style
  requests \\
Architect drawing a typed wiring diagram & 0/6; five exact ties & not run & not recorded &
  Rejected \\
\bottomrule
\end{tabularx}
\end{table}

\para{What V1 shipped} Single checkers that held at parity on held-out tasks did ship, one gate at
a time: acceptance (team 6.0), symptom for bug reports (6.1), repository conventions (6.2), a
settle step in the checker's hand-back (6.3; 17/48 against 10/48 on its targets) and a combination
tester (6.4). Every heavier variant failed its held-out check, including a two-checker version
(3/9 against 5/9) and a settler node (20/40 against 21/40 pooled). A census of 61 misses in 120
runs on 18 October found 33 misreadings of the request, 16 self-checks gentler than the hidden
tests, and 12 gaps in domain knowledge.

\para{The lone lead won on Terminal-Bench too} In V8, team 7.0 (a single lead that works directly
and delegates only when asked) at max effort beat the earlier lead at medium effort on the unread
30 by $+3.17$ (standard error 1.39), with 10 tasks up, 3 down and no always-passed task lost, at
twice the cost. This compares configurations: the lead's profile and its effort changed together,
and that pair is what we ship (effort alone is discussed in \cref{sec:judgment}). Proactive
delegation through Codex's own collaboration protocol, in which the lead spawns helper agents
without being asked, reached 2 of 24 target trials, lost two guard tasks and cost two to four
times as much; its passive form, delegation only on request, is the one in team 7.0. A checking
team tied the lone lead at twice the cost.

\para{Same-model second opinions did not help out of sample} Readers that studied the task before
the work scored $-0.8$ ($\pm 0.7$) pooled on the unread set. Choosers between two finished builds,
GPT or Claude, were right about 55\% of the time: almost always when the failing build had a
visible defect (Claude Opus 22 of 23) and rarely on judgment points (5 of 28). A same-model jury
over 36 pass/fail pairs was right on 12, wrong on 11 and abstained on 13. A
clean-room grader started by the lead passed 3 of 10 target trials against 5 of 10 for the base,
and 6 of the 10 ``all passed'' verdicts it gave were on failing trials. A gated reasoning checker
tied at 12/20.

\para{The harness helps when it answers a state passing runs never show} In headless runs nobody
answers a hand-back. When a run handed back a failure early, the harness answered ``decide
yourself, keep going'' (GO); one task went from 0 of 14 to 6 of 6, and GO shipped. Checks
triggered on every first delivery did not ship. HD on GPT-6.1 gained $+1.54$ ($\pm 1.31$) on the
unread set at 2.09 times the cost, below its pre-registered bar of $+2.0$. O2, the same idea worded
for GPT-6.2, passed 22/28 against 17/28 on the readable tasks that split across draws, where its
second round found real defects, but lost 2.0 on the unread set (15.5 against 17.5, $p = 0.28$),
reworking delivered work in 50 of 60 trials. A global trigger buys its gains inside one failure
family and pays for them on tasks that were already right.

\para{Answer to RQ2} On single-context tasks, teams of the same model added cost but not accuracy.
What helped was narrow: a single gate for specification-style requests, and a runtime answer to a
state (an early failure hand-back) that passing runs never reach.
